\chapter{Geometri segitiga}\label{chap:triangle}

Geometri segitiga adalah kajian tentang sifat-sifat segitiga, termasuk pusat dan lingkaran yang terkait dengannya.

Kita membahas hasil-hasil paling dasar dalam geometri segitiga,
terutama untuk menunjukkan bahwa perangkat yang telah kita kembangkan cukup memadai untuk membuktikan berbagai pernyataan.

\section{Lingkaran luar dan pusat lingkaran luar}

\begin{thm}{Teorema}\label{thm:circumcenter}
Garis-garis sumbu dari sisi-sisi setiap segitiga tak degenerat berpotongan di satu titik.
\end{thm}

Titik perpotongan garis-garis sumbu tersebut disebut \index{pusat lingkaran luar}\emph{pusat lingkaran luar}.
Titik ini merupakan pusat \index{lingkaran luar}\emph{lingkaran luar} segitiga itu;
yaitu, lingkaran yang melalui ketiga titik sudut segitiga.
Pusat lingkaran luar suatu segitiga biasanya dinotasikan dengan~$O$.

\begin{wrapfigure}{o}{29mm}
\centering
\includegraphics{mppics/pic-102}
\end{wrapfigure}


\parit{Bukti.}
Misalkan $\triangle ABC$ tak degenerat.
Misalkan $\ell$ dan $m$ berturut-turut merupakan garis sumbu sisi $[AB]$ dan $[AC]$.

Andaikan $\ell$ berpotongan dengan $m$, katakanlah di $O$.

Mari kita terapkan Teorema~\ref{thm:perp-bisect}.
Karena $O\in\ell$, berlaku $OA\z=OB$, dan karena $O\in m$, berlaku $OA\z=OC$.
Akibatnya, $OB\z=OC$;
yaitu, $O$ terletak pada garis sumbu~$[B C]$.

Tinggal ditunjukkan bahwa $\ell\nparallel m$.
Andaikan sebaliknya, bahwa $\ell\parallel m$.
Karena $\ell\perp(AB)$ dan $m\z\perp(AC)$, Latihan~\ref{ex:perp-perp}\textit{\ref{ex:perp-perp:b}} menyiratkan bahwa
$(AC)\parallel(AB)$, dan menurut Teorema~\ref{thm:parallel}, berlaku $(AC)=(AB)$.
Dengan demikian, $\triangle ABC$ degenerat --- suatu kontradiksi.
\qeds

\begin{thm}[!]{Latihan}\label{ex:unique-cline}
Tunjukkan bahwa terdapat tepat satu lingkaran yang melalui titik-titik sudut suatu segitiga tak degenerat di bidang Euklides.
\end{thm}



\section{Garis tinggi dan ortosenter}

Sebuah \index{garis tinggi}\emph{garis tinggi} segitiga adalah garis yang melalui suatu titik sudut dan tegak lurus terhadap garis yang memuat sisi di hadapannya.
Istilah \index{garis tinggi}\emph{tinggi} juga dapat digunakan untuk jarak dari titik sudut tersebut ke titik kaki tegak lurusnya pada garis yang memuat sisi di hadapannya.

\begin{thm}{Teorema}\label{thm:orthocenter}
Ketiga garis tinggi setiap segitiga tak degenerat berpotongan di satu titik.
\end{thm}

Titik perpotongan garis-garis tinggi disebut \index{ortosenter}\emph{ortosenter};
titik ini biasanya dinotasikan dengan~$H$.

{

\begin{wrapfigure}{o}{34mm}
\vskip-4mm
\centering
\includegraphics{mppics/pic-104}
\end{wrapfigure}

\parit{Bukti.}
Ambil suatu segitiga tak degenerat $A B C$.
Tinjau tiga garis $\ell$, $m$, dan $n$
sedemikian sehingga
\begin{align*}
\ell&\parallel(BC),
&
m&\parallel(CA),
&
n&\parallel(AB),
\\
\ell&\ni A,
&
m&\ni B,
&
n&\ni C.
\end{align*}
Karena $\triangle A B C$ tak degenerat,
tidak ada dua di antara garis $\ell$, $m$, dan $n$ yang sejajar.
Misalkan $A'$, $B'$, dan $C'$ berturut-turut merupakan titik perpotongan
$m$ dengan $n$, $n$ dengan $\ell$, dan $\ell$ dengan $m$.

}

Perhatikan bahwa $\square A B C B'$, $\square B C A C'$, dan $\square C A B A'$ merupakan jajar genjang.
Dengan menerapkan Lema~\ref{lem:parallelogram}, kita memperoleh bahwa $\triangle ABC$ merupakan \index{segitiga medial}\emph{segitiga medial} dari $\triangle A' B' C'$;
yaitu, $A$, $B$, dan $C$ berturut-turut merupakan titik tengah $[B' C']$, $[C' A']$, dan $[A' B']$.

Karena $(B' C')\parallel (BC)$, Latihan~\ref{ex:perp-perp}\textit{\ref{ex:perp-perp:a}} menyiratkan bahwa garis tinggi dari $A$ tegak lurus terhadap $[B' C']$, dan menurut uraian di atas, garis itu membagi dua~$[B' C']$.

Dengan demikian, garis-garis tinggi $\triangle A B C$
juga merupakan garis-garis sumbu $\triangle A' B' C'$.
Dengan menerapkan Teorema~\ref{thm:circumcenter}, kita memperoleh bahwa garis-garis tinggi $\triangle ABC$ berpotongan di satu titik.
\qeds

\begin{thm}{Latihan}\label{ex:orthic-4}
Andaikan ortosenter $H$ dari $\triangle ABC$ berbeda dari titik-titik sudutnya.
Tunjukkan bahwa $A$ merupakan ortosenter $\triangle H B C$.
\end{thm}

\begin{thm}{Latihan}\label{ex:orthic-sim}
Misalkan $A'$, $B'$, dan $C'$ berturut-turut merupakan titik-titik kaki dari ketiga garis tinggi suatu segitiga lancip $ABC$.
Tunjukkan bahwa
\[
\triangle ABC \sim \triangle A'B'C \sim \triangle AB'C' \sim \triangle A'BC'.
\]
\end{thm}



\section{Garis berat dan titik berat}

Garis berat suatu segitiga adalah ruas yang menghubungkan sebuah titik sudut dengan titik tengah sisi di hadapannya.

\begin{thm}{Teorema}\label{thm:centroid}
Ketiga garis berat setiap segitiga tak degenerat berpotongan di satu titik.
Selain itu, titik perpotongan tersebut membagi setiap garis berat dalam perbandingan 2:1.
\end{thm}

Titik perpotongan garis-garis berat disebut \index{titik berat}\emph{titik berat} segitiga;
titik ini biasanya dinotasikan dengan~$M$.
Dalam bukti ini, kita akan menerapkan Latihan \ref{ex:chevinas} dan \ref{ex:smililar+parallel}; penyelesaian lengkap keduanya diberikan dalam petunjuk.

\parit{Bukti.}
Tinjau suatu segitiga tak degenerat $A B C$.
Misalkan $[A A']$ dan $[B B']$ merupakan dua garis beratnya.
Menurut Latihan~\ref{ex:chevinas},
$[A A']$ dan $[B B']$ memiliki suatu titik perpotongan;
nyatakan titik itu dengan $M$.

\begin{wrapfigure}{o}{36mm}
\vskip-4mm
\centering
\includegraphics{mppics/pic-106}
\end{wrapfigure}

Tarik garis $\ell$ yang melalui $A'$ dan sejajar dengan $(BB')$.
Dengan menerapkan Latihan~\ref{ex:smililar+parallel} pada $\triangle BB'C$ dan $\ell$, kita memperoleh bahwa $\ell$ memotong $[B'C]$ di suatu titik, katakanlah $X$, dan
\[\frac{CX}{CB'}=\frac{CA'}{CB}=\frac12;\]
yaitu, $X$ merupakan titik tengah $[CB']$.

Karena $B'$ merupakan titik tengah $[AC]$ dan $X$ merupakan titik tengah $[B'C]$, kita memperoleh
\[\frac{AB'}{AX}=\frac23.\]

Dengan menerapkan Latihan~\ref{ex:smililar+parallel} pada $\triangle XA'A$ dan garis $(BB')$, kita memperoleh
\[\frac{AM}{AA'}=\frac{AB'}{AX}=\frac23;\eqlbl{eq:2/3}\]
yaitu, $M$ membagi $[AA']$ dalam perbandingan 2:1.

Kesamaan \ref{eq:2/3} menentukan $M$ secara tunggal pada $[AA']$.
Dengan mengulangi argumen yang sama untuk garis-garis berat $[AA']$ dan $[CC']$, kita memperoleh bahwa keduanya juga berpotongan di~$M$,
sehingga hasilnya terbukti.
\qeds


\begin{thm}{Latihan}\label{ex:midle}
Misalkan $\square ABCD$ merupakan segiempat tak degenerat
dan $X$, $Y$, $V$, serta~$W$ berturut-turut merupakan titik tengah
$[AB]$, $[BC]$, $[CD]$, dan~$[DA]$.
Tunjukkan bahwa $\square XYVW$ merupakan jajar genjang.
\end{thm}

\begin{thm}{Latihan lanjutan}\label{ex:euler-line}
Tunjukkan bahwa ortosenter $H$, titik berat $M$, dan pusat lingkaran luar $O$ dari setiap segitiga tak degenerat $ABC$ terletak pada satu garis.
Selain itu, titik berat membagi ruas $[HO]$ dalam perbandingan $2:1$.
\end{thm}

Garis dalam latihan ini disebut \emph{garis Euler}.


\section{Garis bagi sudut}

Jika $\measuredangle A B X\equiv-\measuredangle C B X$,
maka kita katakan bahwa garis $(BX)$ {}\emph{membagi dua} $\angle ABC$,
atau bahwa garis $(BX)$ merupakan \index{garis bagi!garis bagi sudut}\emph{garis bagi} $\angle ABC$.
Jika $\measuredangle A B X\equiv\pi-\measuredangle C B X$, maka garis $(BX)$ disebut \index{garis bagi!garis bagi luar}\emph{garis bagi luar} $\angle ABC$.


\begin{wrapfigure}{o}{42mm}
\centering
\includegraphics{mppics/pic-108}
\end{wrapfigure}

Jika $\measuredangle ABA'=\pi$;
yaitu, jika $B$ terletak di antara $A$ dan $A'$,
maka garis bagi $\angle ABC$ merupakan garis bagi luar $\angle A' B C$, dan sebaliknya.

Perhatikan bahwa garis bagi dan garis bagi luar ditentukan secara tunggal oleh sudutnya.

\begin{thm}{Latihan kelas}\label{ex:perp-bisectors}
Tunjukkan bahwa untuk setiap sudut, garis bagi dan garis bagi luarnya saling tegak lurus.
\end{thm}

Garis-garis bagi $\angle ABC$, $\angle BCA$, dan $\angle CAB$ dari suatu segitiga tak degenerat $A B C$
disebut \index{garis bagi!garis bagi segitiga}\emph{garis-garis bagi segitiga} $A B C$ masing-masing di titik sudut $B$, $C$, dan $A$.

\begin{thm}{Latihan}\label{ex:bisect=altitude}
Andaikan pada salah satu titik sudut suatu segitiga tak degenerat, garis baginya berimpit dengan garis tingginya.
Tunjukkan bahwa segitiga tersebut sama kaki.
\end{thm}

\begin{thm}{Lema}\label{lem:bisect-ratio}
Misalkan $\triangle A B C$ merupakan segitiga tak degenerat.
Andaikan garis bagi di $A$
memotong $[BC]$ di~$D$.
Maka
$$\frac{AB}{AC}=\frac{DB}{DC}.
\eqlbl{bisect-ratio}$$

\end{thm}

\begin{wrapfigure}{r}{28mm}
\vskip-6mm
\centering
\includegraphics{mppics/pic-110}
\end{wrapfigure}

\parit{Bukti.}
Misalkan $\ell$ merupakan garis yang melalui $C$ dan sejajar dengan~$(AB)$.
Perhatikan bahwa garis $\ell$ dan $(AD)$ tidak sejajar;
misalkan $E$ merupakan titik perpotongannya.

Perhatikan pula bahwa $B$ dan $C$ terletak pada sisi-sisi yang berlawanan dari~$(AD)$.
Menurut sifat transversal (\ref{thm:parallel-2}),
$$\measuredangle BAD=\measuredangle CED.\eqlbl{eq:<BAD=<CED}$$

Selanjutnya, sudut $ADB$ dan $EDC$ merupakan sudut bertolak belakang; menurut \ref{prop:vert}, berlaku
$$\measuredangle ADB=\measuredangle EDC.$$

Menurut kriteria kesebangunan AA,
$\triangle ABD\sim \triangle ECD$.
Khususnya,
$$\frac{AB}{EC}=\frac{DB}{DC}.\eqlbl{eq:AB/EC=DB/DC}$$

Karena $(AD)$ membagi dua $\angle BAC$, kita memperoleh
$\measuredangle BAD=\measuredangle DAC$.
Bersama dengan \ref{eq:<BAD=<CED},
hal ini menyiratkan bahwa
$\measuredangle CEA=\measuredangle EAC$.
Menurut Teorema~\ref{thm:isos}, $\triangle ACE$ sama kaki;
yaitu, $$EC=AC.$$
Bersama dengan \ref{eq:AB/EC=DB/DC}, hal ini menyiratkan \ref{bisect-ratio}.
\qeds

\begin{thm}{Latihan}\label{ex:ext-disect}
Rumuskan dan buktikan analog dari Lema~\ref{lem:bisect-ratio} untuk garis bagi luar.
\end{thm}

\begin{thm}{Latihan}\label{ex:bisect=median}
Andaikan garis bagi sudut suatu segitiga tak degenerat membagi sisi di hadapannya menjadi dua sama panjang.
Tunjukkan bahwa segitiga tersebut sama kaki.
\end{thm}

{

\begin{wrapfigure}{r}{38mm}
\vskip-5mm
\centering
\includegraphics{mppics/pic-118}
\end{wrapfigure}

\begin{thm}{Latihan}\label{ex:bisector-parallel}
Andaikan garis bagi sudut di $A$ pada suatu segitiga tak degenerat $ABC$ memotong $[BC]$ di $D$;
garis yang melalui $D$ dan sejajar dengan $(CA)$ memotong $(AB)$ di $E$;
garis yang melalui $E$ dan sejajar dengan $(BC)$ memotong $(AC)$ di $F$.

Tunjukkan bahwa
$AE=FC$.

\end{thm}

}

\section{Sifat berjarak sama}

Ingat bahwa jarak antara suatu garis $\ell$ dan titik $P$ didefinisikan sebagai jarak dari $P$ ke titik kaki tegak lurusnya pada $\ell$; lihat Bagian~\ref{sec:perp<oblique}.

\begin{thm}[\abs]{Proposisi}\label{prop:angle-bisect-dist}
Andaikan $\triangle ABC$ tak degenerat.
Suatu titik $X$ terletak pada garis bagi atau garis bagi luar $\angle ABC$
jika dan hanya jika $X$ berjarak sama dari garis $(AB)$ dan $(BC)$.
\end{thm}


\parit{Bukti.}
Kita boleh mengandaikan bahwa $X$ tidak terletak pada gabungan garis $(AB)$ dan $(BC)$.
Jika tidak, jarak ke salah satu garis bernilai nol;
dalam hal ini, $X=B$ merupakan satu-satunya titik yang berjarak sama dari kedua garis.

Misalkan $Y$ dan $Z$ berturut-turut merupakan cerminan $X$ terhadap $(AB)$ dan $(BC)$.
Perhatikan bahwa
\[Y\ne Z.\]
Jika tidak, kedua garis $(AB)$ dan $(BC)$ merupakan garis sumbu $[XY]$.
Akibatnya, $(AB)=(BC)$, tetapi hal ini mustahil karena $\triangle ABC$ tak degenerat.

Menurut Proposisi~\ref{prop:reflection}, $XB=YB=ZB$.
Perhatikan bahwa $X$ berjarak sama dari $(AB)$ dan $(BC)$ jika dan hanya jika $XY\z=XZ$.


{

\begin{wrapfigure}{r}{24mm}
\centering
\includegraphics{mppics/pic-112}
\end{wrapfigure}

Dengan menerapkan SSS lalu SAS, kita memperoleh
$$\begin{aligned}
XY&=XZ
\\
&\hskip0.4mm\Updownarrow
\\
\triangle XBY&\cong\triangle XBZ
\\
&\hskip0.4mm\Updownarrow
\\
\measuredangle XBY&\equiv \pm \measuredangle XBZ.
\end{aligned}
$$

Karena $Y\ne Z$, kita memperoleh bahwa $\measuredangle XBY\ne \measuredangle XBZ$.
Oleh karena itu, $X$ berjarak sama dari $(AB)$ dan $(BC)$ jika dan hanya jika
\[\measuredangle XBY\equiv -\measuredangle XBZ.
\eqlbl{eq:iff-chain}\]

}

Menurut Proposisi~\ref{prop:reflection}, $A$ terletak pada garis bagi $\angle XBY$,
dan $C$ terletak pada garis bagi $\angle XBZ$; yaitu,
\begin{align*}
2\cdot \measuredangle XBA&\equiv \measuredangle XBY,
&
2\cdot \measuredangle XBC&\equiv \measuredangle XBZ.
\end{align*}
Menurut \ref{eq:iff-chain},
\[2\cdot \measuredangle XBA\equiv -2\cdot \measuredangle XBC.\]
Kesamaan terakhir berarti bahwa
\[
\measuredangle XBA+\measuredangle XBC\equiv 0
\quad
\text{atau}
\quad
\measuredangle XBA+\measuredangle XBC\equiv \pi
\]
--- sehingga hasilnya terbukti.
\qeds

\section{Pusat lingkaran dalam}


\begin{thm}[\abs]{Teorema}\label{thm:incenter}
Garis-garis bagi sudut setiap segitiga tak degenerat berpotongan di satu titik.
\end{thm}

Titik perpotongan garis-garis bagi disebut \index{pusat lingkaran dalam}\emph{pusat lingkaran dalam} segitiga,
dan biasanya dinotasikan dengan~$I$.
Titik $I$ berjarak sama dari setiap sisi.
Khususnya, titik ini merupakan pusat sebuah lingkaran yang menyinggung setiap sisi segitiga.
Lingkaran ini disebut
\index{lingkaran dalam}\emph{lingkaran dalam}, dan jari-jarinya disebut
\index{jari-jari lingkaran dalam}\emph{jari-jari lingkaran dalam} segitiga.

\parit{Bukti.}
Misalkan $\triangle ABC$ merupakan segitiga tak degenerat.

Titik $B$ dan $C$ terletak pada sisi-sisi yang berlawanan dari garis bagi $\angle BAC$.
Oleh karena itu, garis bagi ini memotong $[BC]$ di suatu titik, katakanlah~$A'$.

Secara serupa, terdapat $B'\in[AC]$
sedemikian sehingga $(BB')$ membagi dua $\angle ABC$.

Dengan menerapkan Latihan~\ref{ex:chevinas},
kita memperoleh bahwa $[AA']$ dan $[BB']$ berpotongan.
Misalkan $I$ menyatakan titik perpotongan tersebut.

{

\begin{wrapfigure}{o}{28mm}
\vskip0mm
\centering
\includegraphics{mppics/pic-114}
\end{wrapfigure}

Misalkan $X$, $Y$, dan $Z$ berturut-turut merupakan titik kaki tegak lurus dari $I$ pada $(B C)$, $(C A)$, dan $(A B)$.
Dengan menerapkan Proposisi~\ref{prop:angle-bisect-dist}, kita memperoleh
$$I Y=I Z=I X.$$
Dari proposisi yang sama, kita memperoleh bahwa $I$ terletak pada garis bagi atau garis bagi luar $\angle B C A$.

Garis $(C I)$ memotong $[B B']$;
titik $B$ dan $B'$ terletak pada sisi-sisi yang berlawanan dari~$(C I)$.
Oleh karena itu, sudut $I C B'$ dan $I C B$ memiliki tanda yang berlawanan.
Perhatikan bahwa $\angle I C A=\angle I C B'$.
Dengan demikian, $(C I)$ tidak mungkin merupakan garis bagi luar $\angle B C A$.
Maka hasilnya terbukti.
\qeds

}

{

\begin{wrapfigure}{r}{30mm}
\centering
\vskip-2mm
\includegraphics{mppics/pic-116}
\end{wrapfigure}

\begin{thm}{Latihan}\label{ex:2x=b+c-a}
Andaikan sisi-sisi $[B C]$, $[C A]$, dan $[A B]$ dari $\triangle A B C$ menyinggung lingkaran dalam berturut-turut di $X$, $Y$, dan $Z$.
Tunjukkan bahwa
$$AY=AZ= \tfrac12\cdot(A B+ A C- B C).$$

\end{thm}

Andaikan $A'$, $B'$, dan $C'$ merupakan titik-titik kaki dari ketiga garis tinggi suatu segitiga $ABC$ dan semuanya berbeda satu sama lain.
Maka $\triangle A'B'C'$ disebut \index{segitiga!segitiga ortik}\index{segitiga ortik}\emph{segitiga ortik} dari $\triangle ABC$.

}

\begin{thm}{Latihan}\label{ex:orthic-triangle}
Buktikan bahwa ortosenter suatu segitiga lancip berimpit dengan pusat lingkaran dalam segitiga ortiknya.

Apa analog dari pernyataan ini untuk segitiga tumpul?
\end{thm}

\begin{thm}{Latihan}\label{ex:bisector-incenter}
Misalkan $I$ merupakan titik perpotongan garis-garis bagi di $A$ dan $B$ dari suatu segitiga tak degenerat $ABC$.
Misalkan $D$ menyatakan titik perpotongan garis bagi di $A$ dengan sisi $[BC]$.
Tunjukkan bahwa $\frac{AI}{DI}=\frac{b+c}{a}$,
dengan $a=BC$, $b=CA$, dan $c=AB$.

Gunakan hasil ini untuk menyusun bukti lain bagi \ref{thm:incenter}.
\end{thm}
